\documentclass[12pt,a4paper]{article}
\usepackage{fontspec}
\setmainfont{Times New Roman}[Path=C:/Windows/Fonts/,UprightFont=times.ttf,BoldFont=timesbd.ttf,ItalicFont=timesi.ttf,BoldItalicFont=timesbi.ttf]
\usepackage{amsmath,amssymb,amsfonts}
\usepackage[left=1.5cm,right=1.5cm,top=1.5cm,bottom=1.7cm,headheight=17pt]{geometry}
\usepackage{enumitem,array,fancyhdr,booktabs}

\pagestyle{fancy}
\fancyhf{}
\lhead{\small\textbf{ĐÁP ÁN CHUYÊN ĐỀ CHƯƠNG III -- TỨ PHÂN VỊ \& PHƯƠNG SAI}}
\rhead{\small\textit{Biên soạn: Dương Gia Bảo}}
\rfoot{\small Trang \thepage}
\renewcommand{\headrulewidth}{0.4pt}
\fancypagestyle{firstpage}{\fancyhead{}\renewcommand{\headrulewidth}{0pt}}

\setlength{\parindent}{0pt}
\setlength{\parskip}{4pt}
\setlength{\emergencystretch}{1em}

\setlist[enumerate,1]{label=\textbf{Câu \arabic*:},leftmargin=*,itemsep=8pt,topsep=5pt}

\begin{document}
\thispagestyle{firstpage}

\noindent
\begin{minipage}[t]{0.36\textwidth}
    \centering
    {\fontsize{12}{14}\selectfont\textbf{BÀI TẬP CHUYÊN ĐỀ}}\\[4pt]
    \textit{Biên soạn: Dương Gia Bảo}\\[6pt]
    \setlength{\fboxrule}{0.8pt}
    \setlength{\fboxsep}{5pt}
    \fbox{\begin{minipage}{0.88\linewidth}
        \centering
        {\fontsize{12}{14}\selectfont\textbf{ĐÁP ÁN CHƯƠNG III}}
    \end{minipage}}
\end{minipage}\hfill
\begin{minipage}[t]{0.62\textwidth}
    \centering
    {\fontsize{13}{15}\selectfont\textbf{CHUYÊN ĐỀ ÔN TẬP CHƯƠNG III}}\\[3pt]
    {\fontsize{12}{14}\selectfont\textbf{TỨ PHÂN VỊ VÀ PHƯƠNG SAI MẪU SỐ LIỆU GHÉP NHÓM}}\\[3pt]
    {\fontsize{11}{13}\selectfont\textit{(Chương trình GDPT 2018 -- Gồm 30 câu trọng tâm)}}
\end{minipage}

\par\vspace{4pt}
\hrule height 0.6pt
\vspace{8pt}

\begin{center}
    {\fontsize{13}{15}\selectfont\textbf{BẢNG ĐÁP ÁN NHANH 30 CÂU}}
\end{center}

\begin{center}
\renewcommand{\arraystretch}{1.25}
\begin{tabular}{|c|c||c|c||c|c|}
\hline
\textbf{Câu} & \textbf{Đáp án} & \textbf{Câu} & \textbf{Đáp án} & \textbf{Câu} & \textbf{Đáp án} \\ \hline
\textbf{1} & \textbf{B} & \textbf{11} & $32{,}78$ & \textbf{21} & $0{,}44$ \\ \hline
\textbf{2} & \textbf{C} & \textbf{12} & $10$ & \textbf{22} & $0{,}53$ \\ \hline
\textbf{3} & \textbf{C} & \textbf{13} & $167{,}5$ & \textbf{23} & $11{,}42$ \\ \hline
\textbf{4} & \textbf{A} & \textbf{14} & $1{,}44$ & \textbf{24} & $47{,}19$ \\ \hline
\textbf{5} & \textbf{D} & \textbf{15} & $16{,}25$ & \textbf{25} & $4{,}76$ \\ \hline
\textbf{6} & \textbf{A} & \textbf{16} & $62{,}5$ & \textbf{26} & $4$ \\ \hline
\textbf{7} & \textbf{B} & \textbf{17} & $4$ & \textbf{27} & $35100$ \\ \hline
\textbf{8} & \textbf{D} & \textbf{18} & $4{,}64$ & \textbf{28} & $324$ \\ \hline
\textbf{9} & \textbf{A} & \textbf{19} & $14{,}11$ & \textbf{29} & $10$ \\ \hline
\textbf{10} & \textbf{C} & \textbf{20} & $159{,}83$ & \textbf{30} & $56$ \\ \hline
\end{tabular}
\end{center}

\vspace{10pt}
\hrule height 0.6pt
\vspace{10pt}

\begin{center}
    {\fontsize{13}{15}\selectfont\textbf{HƯỚNG DẪN GIẢI CHI TIẾT}}
\end{center}

\subsection*{PHẦN I. TRẮC NGHIỆM NHIỀU LỰA CHỌN (CÂU 1 -- 10)}

\begin{enumerate}
\item \textbf{Đáp án B.}  
Tứ phân vị thứ hai $Q_2$ chính là trung vị $M_e$ của mẫu số liệu ghép nhóm.

\item \textbf{Đáp án C.}  
Theo quy tắc phát hiện giá trị ngoại lệ trong SGK Toán 12, giá trị $x_0$ là ngoại lệ nếu $x_0 > Q_3 + 1{,}5\Delta_Q$ hoặc $x_0 < Q_1 - 1{,}5\Delta_Q$.

\item \textbf{Đáp án C.}  
Cỡ mẫu $n = 40 \implies \dfrac{3n}{4} = 30$.  
Tần số tích lũy: $cf_1 = 5; cf_2 = 5 + 15 = 20 < 30; cf_3 = 20 + 12 = 32 \ge 30$.  
Vậy nhóm chứa $Q_3$ là $[20; 30)$.

\item \textbf{Đáp án A.}  
Cỡ mẫu $n = 20$.  
- $\dfrac{n}{4} = 5 \implies$ nhóm chứa $Q_1$ là $[4; 6)$: $Q_1 = 4 + \dfrac{5 - 4}{12} \cdot 2 = 4 + \dfrac{1}{6} = \dfrac{25}{6}$.  
- $\dfrac{3n}{4} = 15 \implies$ nhóm chứa $Q_3$ là $[4; 6)$: $Q_3 = 4 + \dfrac{15 - 4}{12} \cdot 2 = 4 + \dfrac{11}{6} = \dfrac{35}{6}$.  
Suy ra $\Delta_Q = Q_3 - Q_1 = \dfrac{35}{6} - \dfrac{25}{6} = \dfrac{10}{6} = \dfrac{5}{3}$.

\item \textbf{Đáp án D.}  
Khoảng tứ phân vị $\Delta_Q = Q_3 - Q_1$ đo độ phân tán của $50\%$ số liệu ở giữa và không bị ảnh hưởng bởi các giá trị ngoại lệ.

\item \textbf{Đáp án A.}  
Công thức tính nhanh phương sai của mẫu số liệu ghép nhóm: $s^2 = \dfrac{1}{n}\sum\limits_{i=1}^k n_i c_i^2 - \bar{x}^2$.

\item \textbf{Đáp án B.}  
Phương sai là bình phương của độ lệch chuẩn: $s^2 = (0{,}8)^2 = 0{,}64$.

\item \textbf{Đáp án D.}  
Giá trị đại diện: $2; 4; 6$. Cỡ mẫu $n = 20$. Số trung bình $\bar{x} = \dfrac{5(2) + 10(4) + 5(6)}{20} = 4$.  
Phương sai: $s^2 = \dfrac{5(2-4)^2 + 10(4-4)^2 + 5(6-4)^2}{20} = \dfrac{20 + 0 + 20}{20} = 2$.

\item \textbf{Đáp án A.}  
Ta có $s_B^2 = (2{,}4)^2 = 5{,}76$. Vì $s_A^2 = 4{,}5 < 5{,}76 = s_B^2$ nên mẫu $A$ có độ phân tán nhỏ hơn và đồng đều hơn mẫu $B$.

\item \textbf{Đáp án C.}  
Khi cộng thêm cùng một hằng số $c$ vào tất cả các giá trị (hay các đầu mút của từng nhóm), số trung bình tăng thêm $c$ nhưng độ lệch $(c_i - \bar{x})$ không đổi, do đó phương sai $s^2$ không đổi.
\end{enumerate}

\subsection*{PHẦN II. BÀI TẬP TRẢ LỜI NGẮN (CÂU 11 -- 30)}

\begin{enumerate}[start=11]
\item \textbf{Đáp số: $32{,}78$.}  
Cỡ mẫu $n = 5 + 9 + 13 + 10 + 6 = 43 \implies \dfrac{n}{4} = 10{,}75$.  
Tần số tích lũy: $cf_1 = 5 < 10{,}75$; $cf_2 = 5 + 9 = 14 \ge 10{,}75$.  
Nhóm chứa $Q_1$ là $[20; 40)$:  
\[
Q_1 = 20 + \dfrac{10{,}75 - 5}{9} \cdot 20 = 20 + \dfrac{115}{9} = \dfrac{295}{9} \approx 32{,}78.
\]

\item \textbf{Đáp số: $10$.}  
Cỡ mẫu $n = 100$.  
- Tìm $Q_2$: $\dfrac{n}{2} = 50$. Tần số tích lũy: $cf_2 = 15 + 25 = 40 < 50$; $cf_3 = 40 + 30 = 70 \ge 50$.  
Nhóm chứa $Q_2$ là $[70; 80) \implies Q_2 = 70 + \dfrac{50 - 40}{30} \cdot 10 = 70 + \dfrac{10}{3}$.  
- Tìm $Q_3$: $\dfrac{3n}{4} = 75$. Tần số tích lũy: $cf_4 = 70 + 15 = 85 \ge 75$.  
Nhóm chứa $Q_3$ là $[80; 90) \implies Q_3 = 80 + \dfrac{75 - 70}{15} \cdot 10 = 80 + \dfrac{10}{3}$.  
Vậy $Q_3 - Q_2 = \left(80 + \dfrac{10}{3}\right) - \left(70 + \dfrac{10}{3}\right) = 10$.

\item \textbf{Đáp số: $167{,}5$.}  
Cỡ mẫu $n = 210$.  
- Tìm $Q_1$: $\dfrac{n}{4} = 52{,}5 < 60 \implies$ nhóm chứa $Q_1$ là $[100; 200)$:  
\[
Q_1 = 100 + \dfrac{52{,}5 - 0}{60} \cdot 100 = 187{,}5.
\]
- Tìm $Q_3$: $\dfrac{3n}{4} = 157{,}5$. Tần số tích lũy: $cf_2 = 130 < 157{,}5$; $cf_3 = 180 \ge 157{,}5 \implies$ nhóm chứa $Q_3$ là $[300; 400)$:  
\[
Q_3 = 300 + \dfrac{157{,}5 - 130}{50} \cdot 100 = 355.
\]
Vậy $\Delta_Q = 355 - 187{,}5 = 167{,}5$.

\item \textbf{Đáp số: $1{,}44$.}  
Cả hai lớp đều có cỡ mẫu $n = 26 \implies \dfrac{n}{4} = 6{,}5$ và $\dfrac{3n}{4} = 19{,}5$.  
- \textbf{Lớp 12C:} $cf = [2; 9; 21; 25; 26]$.  
  $Q_1(12\text{C}) = 160 + \dfrac{6{,}5 - 2}{7} \cdot 5 = \dfrac{1142{,}5}{7} \approx 163{,}2143$.  
  $Q_3(12\text{C}) = 165 + \dfrac{19{,}5 - 9}{12} \cdot 5 = 169{,}375$.  
  $\implies \Delta_Q(12\text{C}) \approx 169{,}375 - 163{,}2143 = 6{,}1607$.  
- \textbf{Lớp 12D:} $cf = [5; 14; 22; 24; 26]$.  
  $Q_1(12\text{D}) = 160 + \dfrac{6{,}5 - 5}{9} \cdot 5 = 160 + \dfrac{5}{6} \approx 160{,}8333$.  
  $Q_3(12\text{D}) = 165 + \dfrac{19{,}5 - 14}{8} \cdot 5 = 168{,}4375$.  
  $\implies \Delta_Q(12\text{D}) \approx 168{,}4375 - 160{,}8333 = 7{,}6042$.  
Hiệu: $\Delta_Q(12\text{D}) - \Delta_Q(12\text{C}) \approx 7{,}6042 - 6{,}1607 \approx 1{,}44$.

\item \textbf{Đáp số: $16{,}25$.}  
Cỡ mẫu $n = 40$. Chú ý độ dài các khoảng không đều ($h_1=5, h_2=10, h_3=15, h_4=30$).  
- $\dfrac{n}{4} = 10 \implies$ nhóm chứa $Q_1$ là $[5; 15)$ có $h_2 = 10$: $Q_1 = 5 + \dfrac{10 - 8}{16} \cdot 10 = 6{,}25$.  
- $\dfrac{3n}{4} = 30 \implies$ nhóm chứa $Q_3$ là $[15; 30)$ có $h_3 = 15$: $Q_3 = 15 + \dfrac{30 - 24}{12} \cdot 15 = 22{,}5$.  
Vậy $\Delta_Q = 22{,}5 - 6{,}25 = 16{,}25$.

\item \textbf{Đáp số: $62{,}5$.}  
Cỡ mẫu $n = 40$.  
- $\dfrac{n}{4} = 10 \implies Q_1 = 45 + \dfrac{10 - 4}{16} \cdot 5 = 46{,}875$.  
- $\dfrac{3n}{4} = 30 \implies Q_3 = 50 + \dfrac{30 - 20}{16} \cdot 5 = 53{,}125$.  
Khoảng tứ phân vị $\Delta_Q = 53{,}125 - 46{,}875 = 6{,}25$.  
Ngưỡng ngoại lệ trên: $M = Q_3 + 1{,}5\Delta_Q = 53{,}125 + 1{,}5(6{,}25) = 53{,}125 + 9{,}375 = 62{,}5\text{ cm}$.

\item \textbf{Đáp số: $4$.}  
Từ cỡ mẫu $n = 40 \implies 6 + x + y + 8 = 40 \iff x + y = 26$.  
Ta có $\dfrac{n}{4} = 10$. Vì $Q_1 = 24 \in [20; 30)$ nên nhóm chứa $Q_1$ là $[20; 30)$:  
\[
20 + \dfrac{10 - 6}{x} \cdot 10 = 24 \iff \dfrac{40}{x} = 4 \iff x = 10.
\]
Suy ra $y = 26 - 10 = 16$. Vậy $2x - y = 2(10) - 16 = 4$.

\item \textbf{Đáp số: $4{,}64$.}  
Coi tổng tần suất là $100$.  
- $Q_1$ ứng với $25\% \implies$ nhóm $[3; 6)$: $Q_1 = 3 + \dfrac{25 - 15}{35} \cdot 3 = 3 + \dfrac{6}{7} \approx 3{,}8571$.  
- $Q_3$ ứng với $75\% \implies$ nhóm $[6; 9)$ (tích lũy đến nhóm 2 là $50\%$, nhóm 3 là $80\%$): $Q_3 = 6 + \dfrac{75 - 50}{30} \cdot 3 = 8{,}5$.  
Vậy $\Delta_Q = 8{,}5 - \left(3 + \dfrac{6}{7}\right) = \dfrac{65}{14} \approx 4{,}64$.

\item \textbf{Đáp số: $14{,}11$.}  
Cỡ mẫu $n = 50$.  
- $\dfrac{n}{4} = 12{,}5 \implies$ nhóm $[30; 40)$: $Q_1 = 30 + \dfrac{12{,}5 - 6}{14} \cdot 10 \approx 34{,}6429$.  
- $\dfrac{3n}{4} = 37{,}5 \implies$ nhóm $[40; 50)$: $Q_3 = 40 + \dfrac{37{,}5 - 20}{20} \cdot 10 = 48{,}75$.  
Khoảng tứ phân vị: $\Delta_Q = 48{,}75 - 34{,}6429 \approx 14{,}11$.

\item \textbf{Đáp số: $159{,}83$.}  
Cỡ mẫu $n = 60$.  
- $\dfrac{n}{4} = 15 \implies Q_1 = 40 + \dfrac{15 - 9}{15} \cdot 10 = 44$.  
- $\dfrac{n}{2} = 30 \implies Q_2 = 50 + \dfrac{30 - 24}{18} \cdot 10 = 50 + \dfrac{10}{3} \approx 53{,}3333$.  
- $\dfrac{3n}{4} = 45 \implies Q_3 = 60 + \dfrac{45 - 42}{12} \cdot 10 = 62{,}5$.  
Tổng $Q_1 + Q_2 + Q_3 = 44 + 53{,}3333 + 62{,}5 = 159{,}8333 \approx 159{,}83$.

\item \textbf{Đáp số: $0{,}44$.}  
Giá trị đại diện: $0{,}25; 0{,}75; 1{,}25; 1{,}75; 2{,}25$. Cỡ mẫu $n = 100$.  
\begin{align*}
\bar{x} &= \dfrac{20(0{,}25) + 15(0{,}75) + 30(1{,}25) + 20(1{,}75) + 15(2{,}25)}{100} = \dfrac{122{,}5}{100} = 1{,}225. \\
s^2 &= \dfrac{20(0{,}25)^2 + 15(0{,}75)^2 + 30(1{,}25)^2 + 20(1{,}75)^2 + 15(2{,}25)^2}{100} - (1{,}225)^2 \\
&= 1{,}9375 - 1{,}500625 \approx 0{,}44.
\end{align*}

\item \textbf{Đáp số: $0{,}53$.}  
Giá trị đại diện: $19{,}25; 19{,}75; 20{,}25; 20{,}75; 21{,}25$. Cỡ mẫu $n = 100$.  
\[
\bar{x} = \dfrac{13(19{,}25) + 45(19{,}75) + 24(20{,}25) + 12(20{,}75) + 6(21{,}25)}{100} = 20{,}015.
\]
Phương sai $s^2 = \dfrac{1}{100}\sum n_i c_i^2 - (20{,}015)^2 = 400{,}8775 - 400{,}600225 = 0{,}277275$.  
Độ lệch chuẩn: $s = \sqrt{0{,}277275} \approx 0{,}5266 \approx 0{,}53$.

\item \textbf{Đáp số: $11{,}42$.}  
Giá trị đại diện: $12{,}5; 17{,}5; 22{,}5; 27{,}5; 32{,}5; 37{,}5$. Cỡ mẫu đều bằng $60$.  
- \textbf{Công ty A:} $\bar{x}_A = \dfrac{1225}{60} \approx 20{,}4167$; $\dfrac{1}{60}\sum n_i c_i^2 = \dfrac{27962{,}5}{60} \approx 466{,}0417$.  
  $\implies s_A^2 \approx 466{,}0417 - (20{,}4167)^2 = 49{,}1389$.  
- \textbf{Công ty B:} $\bar{x}_B = \dfrac{1150}{60} \approx 19{,}1667$; $\dfrac{1}{60}\sum n_i c_i^2 = \dfrac{25675}{60} \approx 427{,}9167$.  
  $\implies s_B^2 \approx 427{,}9167 - (19{,}1667)^2 = 60{,}5556$.  
Hiệu phương sai: $s_B^2 - s_A^2 \approx 60{,}5556 - 49{,}1389 = 11{,}4167 \approx 11{,}42$.

\item \textbf{Đáp số: $47{,}19$.}  
Giá trị đại diện: $152{,}5; 157{,}5; 162{,}5; 167{,}5; 172{,}5; 177{,}5$. Cỡ mẫu $n = 43$.  
\begin{align*}
\bar{x} &= \dfrac{5(152{,}5) + 10(157{,}5) + 12(162{,}5) + 9(167{,}5) + 4(172{,}5) + 3(177{,}5)}{43} \\
&= \dfrac{7017{,}5}{43} \approx 163{,}1977.
\end{align*}
Phương sai: $s^2 = \dfrac{1147268{,}75}{43} - \left(\dfrac{7017{,}5}{43}\right)^2 \approx 26680{,}6686 - 26633{,}4819 \approx 47{,}19$.

\item \textbf{Đáp số: $4{,}76$.}  
Từ cỡ mẫu $n = 40 \implies 2 + x + 16 + y + 6 = 40 \iff x + y = 16$.  
Giá trị đại diện: $3; 5; 7; 9; 11$.  
Từ $\bar{x} = 7{,}2 \implies \dfrac{2(3) + 5x + 16(7) + 9y + 6(11)}{40} = 7{,}2 \iff 184 + 5x + 9y = 288 \iff 5x + 9y = 104$.  
Thay $y = 16 - x$: $5x + 9(16 - x) = 104 \iff -4x = -40 \iff x = 10 \implies y = 6$.  
Phương sai:  
\[
s^2 = \dfrac{2(3-7{,}2)^2 + 10(5-7{,}2)^2 + 16(7-7{,}2)^2 + 6(9-7{,}2)^2 + 6(11-7{,}2)^2}{40} = \dfrac{190{,}4}{40} = 4{,}76.
\]

\item \textbf{Đáp số: $4$.}  
Khi nhân mỗi giá trị dữ liệu lên $k = 2$ lần và cộng thêm $10$, phương sai tăng lên $k^2 = 2^2 = 4$ lần.  
Kiểm chứng trực tiếp: $s_X^2 = 2$; mẫu $Y$ có giá trị đại diện $14; 18; 22$, $\bar{y} = 18$, $s_Y^2 = \dfrac{5(16) + 10(0) + 5(16)}{20} = 8$.  
Vậy $\dfrac{s_Y^2}{s_X^2} = \dfrac{8}{2} = 4$.

\item \textbf{Đáp số: $35100$.}  
Giá trị đại diện các nhóm (chú ý độ dài khoảng không đều): $150; 300; 500; 800$. Cỡ mẫu $n = 50$.  
\[
\bar{x} = \dfrac{10(150) + 20(300) + 15(500) + 5(800)}{50} = \dfrac{19000}{50} = 380.
\]
Phương sai:  
\begin{align*}
s^2 &= \dfrac{10(150 - 380)^2 + 20(300 - 380)^2 + 15(500 - 380)^2 + 5(800 - 380)^2}{50} \\
&= \dfrac{1755000}{50} = 35100.
\end{align*}

\item \textbf{Đáp số: $324$.}  
Giá trị đại diện: $20; 40; 60; 80$. Tần suất $f_1 = 0{,}1; f_2 = 0{,}3; f_3 = 0{,}4; f_4 = 0{,}2$.  
\[
\bar{x} = 0{,}1(20) + 0{,}3(40) + 0{,}4(60) + 0{,}2(80) = 2 + 12 + 24 + 16 = 54.
\]
\begin{align*}
s^2 &= 0{,}1(20 - 54)^2 + 0{,}3(40 - 54)^2 + 0{,}4(60 - 54)^2 + 0{,}2(80 - 54)^2 \\
&= 115{,}6 + 58{,}8 + 14{,}4 + 135{,}2 = 324.
\end{align*}

\item \textbf{Đáp số: $10$.}  
Do tính đối xứng của bảng phân phối, $\bar{x} = 25$. Giá trị đại diện là $15; 25; 35$. Cỡ mẫu $n = 2x + 20$.  
Phương sai:  
\[
s^2 = \dfrac{x(15 - 25)^2 + 20(25 - 25)^2 + x(35 - 25)^2}{2x + 20} = \dfrac{200x}{2x + 20} = \dfrac{100x}{x + 10}.
\]
Theo đề bài: $\dfrac{100x}{x + 10} = 50 \iff 100x = 50x + 500 \iff 50x = 500 \iff x = 10$.

\item \textbf{Đáp số: $56$.}  
Gộp hai tổ ta thu được mẫu chung gồm $n = 50$ nhân viên với tần số tương ứng từng nhóm:  
- Nhóm $[10; 20)$: $n_1 = 4 + 6 = 10$ ($c_1 = 15$).  
- Nhóm $[20; 30)$: $n_2 = 10 + 10 = 20$ ($c_2 = 25$).  
- Nhóm $[30; 40)$: $n_3 = 6 + 14 = 20$ ($c_3 = 35$).  
Số trung bình chung: $\bar{x} = \dfrac{10(15) + 20(25) + 20(35)}{50} = \dfrac{1350}{50} = 27$.  
Phương sai chung:  
\[
s^2 = \dfrac{10(15 - 27)^2 + 20(25 - 27)^2 + 20(35 - 27)^2}{50} = \dfrac{1440 + 80 + 1280}{50} = \dfrac{2800}{50} = 56.
\]
\end{enumerate}

\end{document}
